\documentclass[11pt,a4paper]{article}

\usepackage{kotex}
\usepackage{fontspec}
\setmainfont{Times New Roman}
\setmainhangulfont{AppleMyungjo}
\setsanshangulfont{Apple SD Gothic Neo}

\usepackage[a4paper,top=18mm,bottom=18mm,left=20mm,right=20mm]{geometry}
\usepackage{array}
\usepackage{enumitem}
\linespread{1.18}
\setlength{\parindent}{0pt}

\newcommand{\sol}[2]{\vspace{6pt}{\sffamily\bfseries #1} \quad 정답 \fbox{#2}\par\vspace{1pt}}

\begin{document}

\begin{center}
{\fontsize{15}{17}\selectfont\sffamily\bfseries 2026 고1 영어 변형 — 지문 41-42 정답 및 해설}\\[2pt]
{\sffamily 스포츠의 예측 불가능성}
\end{center}
\vspace{6pt}

\begin{center}
\renewcommand{\arraystretch}{1.4}
\begin{tabular}{|c|c|c|c|c|}
\hline
\sffamily 문항 & 1 & 2 & 3 & 4 \\\hline
\sffamily 유형 & \sffamily 영어제목 & \sffamily 어휘 부적절 & \sffamily 빈칸추론 & \sffamily 서술형 \\\hline
\sffamily 정답 & ③ & ③ & ① & (아래) \\\hline
\sffamily 배점 & 2점 & 3점 & 2점 & 2점 \\\hline
\end{tabular}
\end{center}
\vspace{8pt}\hrule\vspace{4pt}

\sol{1. 영어 제목}{③}
글은 스포츠의 예측 불가능성이 관중을 끌어들이고 흥행과 수익을 높이는 \textbf{축복}인 동시에, 경기 품질의 변동성(들쭉날쭉함)을 낳아 마케터에게 \textbf{도전 과제}가 됨을 함께 다룬다. 즉 예측 불가능성은 양면적 성격을 지닌다.\par
③ \textit{Unpredictability: The Double-Edged Sword of Sport's Appeal}(예측 불가능성: 스포츠 매력의 양날의 검) — `이로움이자 도전'이라는 글의 양면성을 정확히 포착. \textbf{정답}.\par
① 상업적 예측 가능성이 스포츠 경쟁을 \textit{훼손한다}는 단정은 본문에 없다. 본문은 예측 가능성 자체를 비판하지 않는다.\par
② 예측 불가능성의 `숨은 비용'만 강조하여 긍정적 측면(흥행·수익)을 누락한 한쪽 면만의 제목이다.\par
④ 스타 선수는 마케터가 품질을 보완하는 한 가지 예시(세부 사항)일 뿐 글 전체의 주제가 아니다.\par
⑤ 팬과 수익은 글의 결론부 한 요소일 뿐, 글 전체가 `성공의 척도'를 논하는 것은 아니다.

\vspace{4pt}
\sol{2. 어휘 부적절}{③}
③ (c) \textbf{presence} → \textbf{absence}. 바로 앞 문장들은 결과를 \textbf{예측할 수 없을 때}(unpredictability) 관중과 수익이 늘어난다고 했고, 이러한 예측 \textbf{불가능성}이 경기 품질의 변동성으로 이어진다는 흐름이다. 따라서 변동성을 낳는 것은 예측 가능성의 `존재(presence)'가 아니라 `부재(absence)' = 예측 불가능성이어야 논리가 성립한다. \textbf{정답}.\par
① (a) valued: 상업과 달리 스포츠에서는 예측 가능성이 늘 `가치 있게 여겨지는' 것은 아니라는 맥락 → 적절.\par
② (b) attract: 예측 불가능성에 의존해 사람들을 경기로 `끌어들인다'는 맥락 → 적절.\par
④ (d) overcome: 마케터가 품질 변동이라는 문제를 `극복하려' 시도한다는 맥락 → 적절.\par
⑤ (e) financial: 불확실성이 `재정적' 성공에 중요한 요소라는 맥락 → 적절.

\vspace{4pt}
\sol{3. 빈칸추론}{①}
빈칸 문장: \textit{``Sport events actually \underline{\hspace{3cm}} that unpredictability in order to attract people to the game.''}\par
앞 문장에서 스포츠 경험은 예측 불가능할 때 관중에게 더 좋게 경험된다고 했고, 뒤 문장에서는 결과를 예측할 수 없을 때 관중과 수익이 늘어난다고 설명한다. 따라서 스포츠 행사는 사람들을 끌어들이기 위해 그 예측 불가능성에 \textbf{의존한다}는 뜻이 되어야 한다.\par
① \textbf{rely on} — ``depend on''과 같은 의미로 문맥에 가장 적절하다. \textbf{정답}.\par
② eliminate(제거하다), ④ ignore(무시하다)는 글의 핵심과 반대이고, ③ standardize(표준화하다)는 예측 불가능성과 충돌한다. ⑤ predict(예측하다)는 뒤의 \textit{unpredictability}와 의미상 맞지 않는다.

\vspace{4pt}
{\sffamily\bfseries 4. 서술형}\par\vspace{2pt}
\fbox{모범답안} \quad \textbf{스포츠 경기의 불확실한 특성은 전체적인 경험의 한 요소에 불과하지만, 그것은 재정적 성공에 있어 중요한 요소다.}\par\vspace{3pt}
\textbf{채점 기준 (각 요소별 부분 점수):}\par
\begin{itemize}[leftmargin=1.5em,itemsep=1pt,topsep=2pt]
\item \textit{The uncertain nature of sport contests} = `스포츠 경기의 불확실한 특성(본질)'
\item \textit{is just one element of a total experience} = `전체(적인) 경험의 한 요소에 불과하다'
\item \textit{although} = `비록 $\sim$이지만(하지만)' — 양보 접속사를 올바르게 해석
\item \textit{it is an important element for financial success} = `재정적 성공에 있어 중요한 요소다'
\end{itemize}
\vspace{2pt}\textbf{해설:} \textit{although}가 이끄는 양보의 부사절을 정확히 처리하는 것이 핵심이다. 주절은 `불확실성은 한 요소에 불과하다'는 점을, 종속절은 `그럼에도 재정적 성공에 중요하다'는 점을 대비시킨다.

\end{document}
